Los experimentos realizados muestran que existe un claro trade-off entre calidad de solución y eficiencia. Fuerza bruta garantiza el óptimo pero es inviable para instancias medianas o grandes. Programación dinámica también obtiene el óptimo y escala razonablemente, aunque su tiempo crece linealmente con $n$, $M$, $E$ y $q$ simultáneamente, lo que la hace lenta para instancias con recursos máximos. Los algoritmos greedy son extremadamente rápidos en todos los casos, pero sacrifican calidad: greedy1 alcanza solo entre un 22\% y 78\% del óptimo, mientras que greedy2 logra entre 61\% y 98\%, demostrando que un criterio local más informado como $v/(t+c)$ produce soluciones notablemente mejores sin costo computacional adicional significativo.